\section{Explicit calibration and an all-height onset}
\label{sec:calibration}
The constants used in the preceding comparisons can be chosen
explicitly. The point of this section is their uniform domains:
a scalar evaluation at a large height would not prove a disk bound,
an all-order coefficient estimate, or a statement through every change
of \(p\) and \(N\). We give the calibration of \cite{rec11} and its
continuous majorants before using it in the effective construction.

Set
\[
 W=R+Y,\qquad d_{\rm ch}=\lceil W\rceil,\qquad
 \nu=d_{\rm ch}+1\ge4.
\]
The chart integer \(d_{\rm ch}\) is distinct from the displacement
\(d=z-X\). A sufficient rational record is
\begin{equation}
\begin{gathered}
 H=A=C_\Gamma=c_{\mathrm o}=1,\quad D=4,\quad D_0=3,
 \quad g_-=1/64,\\
 C_H=M_H^*=1/4,\quad C_{\rm tail}=640,\quad
 C_\chi=3^{d_{\rm ch}},\quad C_{\rm cal}=32+d_{\rm ch},
 \quad X_{\rm cl}=64\nu.
\end{gathered}
\label{eq:explicit-calibration}
\end{equation}
Here \(M_H^*\) bounds the complete source mass \(M_H\), rather than
denoting its exact value. The domain statements follow.

\subsection{A complex-disk estimate before differentiation}
\label{subsec:disk-calibration}
The digamma remainder \cite[\S5.11(ii)]{dlmfgamma} gives
\begin{equation}
 \psi(t)=\Log t-\frac1{2t}+E(t),\qquad
 \abs{E(t)}\le\frac{\sec^3(\arg t/2)}{12\abs t^2},
 \quad\abs{\arg t}<\pi.
 \label{eq:digamma-calibration}
\end{equation}
For \(t_\pm=(1\pm\ii z)/4\) on \(\abs{z-X}\le X/2\),
one has \(\abs t\ge X/8\), \(\abs{\arg t}<3\pi/4\).
Since \(\sec(3\pi/8)<3\), \(\abs E\le144/X^2\).
Each point of the inner \(X/4\) disk has a \(t\)-disk of radius
\(X/16\) in this larger region. Thus \(\abs{E'}\le2304/X^3\).
Using \(\abs t\ge3X/16\) and
\(\abs w,\abs{w-1}\ge3X/8\) on the inner disk in the exact
formula for \(\ell_\pm''\) proves
\begin{equation}
 \sup_{\abs{z-X}\le X/4}\abs{\ell_\pm''(z)}
 \le\frac1{3X}+\frac{40}{9X^2}+\frac{144}{X^3}
 \le\frac1X\qquad(X\ge32).
 \label{eq:explicit-disk}
\end{equation}
At \(X=32\), multiplication by \(X\) gives \(353/576<1\);
the bound decreases thereafter. Cauchy's estimate now supplies all
orders simultaneously:
\[
 \abs{c_{\epsilon k}}\le
       \frac{4^{k-1}}{k(k+1)}X^{-k}\le4^{k-1}X^{-k}.
\]

At a real center, direct differentiation of
\eqref{eq:provider-logM} yields
\begin{align*}
 \ell_+'(X)&=-\pi/8+\ii L/4+\delta(X),\\
 \delta(X)&=\frac{2X}{1+X^2}+\frac\ii4\Log(1-\ii/X)
             -\frac\ii{8t}+\frac\ii4E(t),\quad t=(1+\ii X)/4.
\end{align*}
Here \(\abs E\le4/X^2\), and for \(X\ge32\),
\begin{equation}
 \abs\delta\le\frac2X+\frac1{4(X-1)}+
                  \frac1{2X}+\frac1{X^2}<\frac3X.
 \label{eq:explicit-slope}
\end{equation}
Reflection gives the other arm. Along real \(u\ge32\),
\(\abs{(\log G)'(u)}\le\pi/8+3/u<1\). Hence
\(G(X+q)\le g\e^{\abs q}\) for real \(\abs q\le X/2\),
\(X\ge64\), proving \(C_\Gamma=1\).

Stirling's formula with its first omitted-term bound gives exactly
\begin{equation}
 G(X)=\frac{\pi^{1/4}}{32}(1+X^2)^{7/8}\e^{-\pi X/8}
  \exp\left\{\frac X4\arctan(1/X)-\frac14+\Rea R_\Gamma(t)\right\},
 \quad\abs{R_\Gamma(t)}\le\frac2{3X}.
 \label{eq:exact-real-gamma-calibration}
\end{equation}
The final exponent exceeds \(-1/2\) for \(X\ge32\), giving
\(g\ge X^{7/4}\e^{-\pi X/8}/64\).
For \(X\ge\max(64,4W)\), integrate \((\log G)'\) from the real
center through the buffered rectangle and use \eqref{eq:explicit-disk}.
Its modulus stays below one, so both
\(g/\abs{G(z)}\) and \(\abs{G(z)}/g\) are at most
\(\e^W\le3^{d_{\rm ch}}\).

\subsection{Bounds for the complete real source}
The beta integral and half-gamma modulus \cite{dlmfbeta,dlmfhalf}
give, for real \(v\),
\[
 \left|\frac{\Gamma(1/4+\ii v)}{\Gamma(1/2+\ii v)}\right|
 \le\frac{\Gamma(1/4)}{\Gamma(1/2)},\qquad
 \abs{\Gamma(1/4+\ii v)}<8\e^{-\pi\abs v/2}.
\]
For the last inequality split the positive integral for \(\Gamma(1/4)\)
at one to obtain \(\Gamma(1/4)<5\), and use
\(\cosh(\pi v)\ge\e^{\pi\abs v}/2\).
The finite-sum zeta identity \cite{dlmfzeta} at \(N=1\) reads
\[
 \zeta(s)=\frac{s}{s-1}-s\int_1^\infty\{x\}x^{-s-1}\dd x.
\]
For \(s=(1+\ii u)/2\), it gives
\(\abs{\zeta(s)}\le1+2\abs s\le2(1+\abs u)\).
The exact normalization \(H(0,u)=M(s)\zeta(s)\) therefore yields
\begin{equation}
 \abs{H(0,u)}\le\tfrac14(1+\abs u)^4\e^{-\pi\abs u/8}
 \qquad(u\in\R).
 \label{eq:global-source-calibration}
\end{equation}
Every summand of \(\Phi(v)\) is positive for \(v\ge0\), since
\(2\pi n^2\e^{4v}>3\). Thus, without a source cutoff,
\begin{equation}
 M_H=\int_0^\infty\abs\Phi=H(0,0)\le1/4.
 \label{eq:complete-source-mass}
\end{equation}
For \(\abs q\le X/2\), \(X\ge32\),
\(1+X+q\le(5/2)X\). Dividing
\eqref{eq:global-source-calibration} by the lower gamma scale gives
\[
 \frac{\abs{H(0,X+q)}}g
     \le640X^{9/4}\e^{-\pi q/8}.
\]
Consequently the moderate and remote constants can be taken as
\begin{equation}
 A_{\rm mod}^{\theta}=640\max(\mathsf D_0^\theta,
                    (\mathsf D_0^\theta)^t),\qquad
 A_{\rm far}=16\max(\mathsf J_0,\mathsf J_0^t).
 \label{eq:explicit-tail-calibration}
\end{equation}

\subsection{Calibration of the actual outer envelope}
Put \(L_z=\Log(z/(4\pi))\),
\(G_0(z)=\pi^{1/4}z^{7/4}\e^{-\pi z/8}/32\), and
\(\theta_0=z(L_z-1)/4+7\pi/8\).
On the buffered rectangle, when \(X\ge\max(16,2(W+2))\),
the exact sectorial formula gives
\begin{equation}
 \ell_\pm(z)=\log G_0(z)\pm\ii\theta_0(z)+\delta_\pm(z),
 \qquad\abs{\delta_\pm(z)}\le8/\abs z\le16/X.
 \label{eq:complex-gamma-calibration}
\end{equation}
For the plus arm the remainder is
\[
 \Log(1+z^{-2})+\frac{\ii z-1}{4}\Log(1-\ii/z)
                      -\frac14+R_\Gamma(t_+).
\]
The logarithm series and \(\abs{R_\Gamma(t_+)}\le16/(3\abs z)\)
bound this by \(8/\abs z\); conjugate reflection gives the minus arm.
The branches are fixed at positive real height. Averaging preserves the
canonical \(G\) and gives \(\abs{\log(G/G_0)}\le16/X\).

The original amplitude is
\(A_{\mathrm o}(-T,z)=2G_0(z)\e^{-TL_z^2/16+\pi^2T/64}\).
Since \(\abs{L_z-L}\le2W/X\),
\(\abs{L_z^2-L^2}\le4W(1+L)/X\) for \(L\ge0\), one obtains
\begin{equation}
 \left|\log\frac{\abs{A_{\mathrm o}(-T,z)}}
 {2\abs{G(z)}\e^{-TL^2/16+\pi^2T/64}}\right|
 \le\frac{C_{\rm cal}}2\frac{1+L}{X}.
 \label{eq:explicit-envelope-calibration}
\end{equation}
Under \(X\ge X_{\rm cl}\) and
\(C_{\rm cal}(1+L)/X\le1\), AM--GM for the two original arms gives
\begin{equation}
 \calE_{\mathrm o}(-T,z)\ge\abs{A_{\mathrm o}(-T,z)}
 \ge\abs{G(z)}\e^{-TL^2/16+\pi^2T/64}.
 \label{eq:explicit-envelope-lower}
\end{equation}
The amplitude ratio in \eqref{eq:explicit-envelope-calibration}
differs from one by at most \(C_{\rm cal}(1+L)/X\).
This proves \(c_{\mathrm o}=1\) while retaining the transverse growth
of the actual two-arm envelope. The common center onset \(64\nu\)
covers the full gamma disk, real window, and buffered chart;
the remaining height-dependent structural gates are paid next.

\subsection{Continuous majorants through every order jump}
\label{subsec:explicit-onset}
For clarity first retain a valid general calibration record.
Let \(\theta=1/2\), \(C_{\mathrm o}=C_\chi/c_{\mathrm o}\),
and choose fixed constants \(S,M\) so that
\begin{align}
 \log\mathcal J(p),\log\mathcal J_\theta(p)
   &\le S(p+1)\log(p+2),\qquad S\ge1,\nonumber\\
 M&=\max\{0,\log^+(8C_{\mathrm o}C_Q),
       \log^+(392C_{\mathrm o}C_{\rm pop}),\nonumber\\
 &\hspace{36mm}\log^+(392C_{\mathrm o}A_{\rm mod}^\theta),
       \log^+(392C_{\mathrm o}A_{\rm far})\},
 \label{eq:onset-log-costs}
\end{align}
where \(Q_p\le C_Q\mathcal J(p)\) and
\(A_{\rm art,p}^\theta\le C_{\rm pop}\mathcal J_\theta(p)\).
For example \(C_Q=48M_0+M_0^t+49C_{\rm pop}\e\) works.
The factorial estimates in \cref{sec:growing} supply these constants.
With \(A=\max(1,H),D=4\), a sufficient general onset is
\begin{align}
 L_{\rm rate}&=\max\{16,(32S)^2,\sqrt{32M}\},\nonumber\\
 L_0&=\max\bigl\{L_{\rm rate},(2R+W+1)/8,\log(1+4(W+1)),
       2\log(1+24A),\nonumber\\
 &\hspace{22mm}2\log(1+D_0),2\log(1+C_{\rm cal}),
                  \log(1+X_{\rm cl}/(4\pi))\bigr\}.
 \label{eq:parameterized-onset}
\end{align}
Here \(X_{\rm cl}\) must already be a valid \emph{center} onset
for every classical domain, including the real window.

Indeed, for \(p=\lfloor L/2\rfloor\), \(L\ge16\),
\[
 (p+1)\log(p+2)\le L\log(L+2)\le L^{3/2},\qquad
 (p+3/4)\log X-3L^2/8\ge7L^2/64.
\]
At \(L\ge L_{\rm rate}\), \(M+SL^{3/2}\le L^2/16\),
so the local half-target gate holds strictly.
The moderate and artificial gates have available quadratic coefficient
\(61/8\); their height powers and logarithmic costs fit with slack.
For the remote gate, \(\pi X/8\le X^2/512\) and
\(X^2/512\ge9L^2/8>(3/8+1/32)L^2\).
These are the four conditions \eqref{eq:growing-gates} with factors
\(8,392\), reserving half the total target for evaluation.

The structural conditions also follow from \eqref{eq:parameterized-onset}.
For \(L\ge16\), \(V\le12L\), \(X\ge12\e^L\), and
\(\sqrt{2\pi X}\ge8\e^{L/2}\). Hence
\[
 \frac{4(V+2)}{\sqrt{2\pi X}}\le(6L+1)\e^{-L/2}<\frac12,
 \quad2(V+2)\le X/4,\quad K\le L\le(X-V)/(4\pi).
\]
The buffer terms in \(L_0\) give the window gates;
\(L\ge2\log(1+24A)\) pays \(2AV^2/X\le1\).
The slope terms give \(X\ge10D_0\) and \(X\ge2D_0(L+2)\),
which pay both the stronger real-slope gate and the separate damping
condition in \eqref{eq:structural-gates}.
The calibration term pays \(C_{\rm cal}(1+L)/X\le1\).
Strong geometry gives the remaining ellipse/logarithm bounds.
No monotonicity of a factorial expression across floor jumps is assumed.

\begin{theorem}[Explicit sufficient onset]\label{thm:explicit-onset}
For the record \eqref{eq:explicit-calibration}, let
\begin{equation}
 L_*=2^{22}\nu^4,\qquad X_*=2^{2L_*+8}.
 \label{eq:closed-onset}
\end{equation}
Every center with \(L\ge L_*\) meets the uniform classical domains,
structural gates, and half-target truncation gates for
\(p=\lfloor L/2\rfloor,K=2p,1\le T\le2\).
The condition \(X\ge X_*\) is a stronger explicit dyadic sufficient
condition. Finite evaluation additionally retains
\eqref{eq:descriptor-certificates} and its separation dependence.
\end{theorem}
\begin{proof}
Use the rational majorants
\[
 \overline E=3^{8\nu^2},\quad\overline S=32\nu,\quad
 \overline Q=32\overline S^2=2^{15}\nu^2,\quad
 \overline T=\overline E\,\overline S^2,
\]
and \(\overline{\mathcal J}(p)=
\overline T(p+2)\overline Q^{p+1}(2p+4)!\).
They dominate both moment systems in \cref{sec:growing}:
\(1+R/4\le\nu/2\), the largest refined-tail exponent is at most
\(9\nu^2/4\), and \(\overline S\) dominates every moment base.
Also the RS constants satisfy \(Q_R<64,C_R<32\).
The factorial bound gives
\[
 \log\overline{\mathcal J}(p)
 \le(6+\log\overline Q+\log\overline T)(p+1)\log(p+2),
 \quad\log\overline T\le18\nu^2,\quad\log\overline Q\le2\nu^2.
\]
The explicit moments imply
\(M_0\le\overline E\), \(M_0^t\le\overline T\),
\(\mathsf D_0^\theta,(\mathsf D_0^\theta)^t,
\mathsf J_0,\mathsf J_0^t\le\overline T\).
For the remote moments one may use
\(\mathsf J_m\le2\e^{Y^2/4}m!(Y+4)^m\).
It follows that
\[
 C_Q\le48\overline E+\overline T+588\le3\overline T,
 \quad A_{\rm mod}^\theta\le640\overline T,
 \quad A_{\rm far}\le16\overline T.
\]
Since \(C_{\mathrm o}=3^{d_{\rm ch}}\), every fixed logarithmic
cost in \eqref{eq:onset-log-costs} is bounded using a multiplier
below \(2^{18}\). Thus \(S,M\le32\nu^2\), and
\(L_{\rm rate}\le2^{20}\nu^4\).
All remaining terms of \eqref{eq:parameterized-onset} are smaller than
\(2^{22}\nu^4\). Finally \(\log2>1/2\), \(\log(4\pi)<3\)
give \(\log(X_*/(4\pi))>L_*+1\).
\end{proof}
The threshold is deliberately conservative. Its role is effectiveness,
not a practical starting height. The radii and any inherited regulator
dependence are fixed before height grows. Neither an endpoint-uniform
time theorem nor a uniform lower bound on center separation follows.

At a smaller specified height there is also a finite verifier: recompute
outward rational enclosures for \eqref{eq:descriptor-certificates}, all
structural/domain gates, and the four half-target inequalities, and
accept only certified strict margins. A finite unsuccessful witness is
rejected or inconclusive. Every strictly admissible height has a finite
successful witness. This does not promote a single-height verification
to an all-height theorem.

\subsection{A smaller onset from additive costs}
\label{subsec:linear-onset}
The preceding onset is a valid fallback. Its main loss comes from
multiplying the fixed chart cost \(\log\overline T\) by
\((p+1)\log(p+2)\), followed by \(L\log(L+2)\le L^{3/2}\).
Keeping that cost additive gives the following improvement
\cite{rec16}. The moment and factorial majorants themselves need not change.

\begin{theorem}[Linear sufficient log-height]\label{thm:linear-onset}
With the same calibration record and frozen choices, every center with
\begin{equation}
 L\ge L_{\rm new}=64\nu
 \label{eq:linear-onset}
\end{equation}
meets all classical, structural, calibration and half-target truncation
gates of \cref{thm:explicit-onset}. A stronger dyadic sufficient condition is
\begin{equation}
 X\ge X_{\rm new}=2^{128\nu+8}.
 \label{eq:linear-dyadic-onset}
\end{equation}
All input and descriptor restrictions remain. In particular the theorem
does not admit excluded boundary centers or bound their separation uniformly.
\end{theorem}
\begin{proof}
Keep \(\overline E,\overline S,\overline Q,\overline T\) and
\(\overline{\mathcal J}(p)\) from the preceding proof.
They dominate both moment systems, with
\[
 Q_p\le C_Q\overline{\mathcal J}(p),\quad C_Q\le3\overline T,
 \quad A_{\rm art,p}^{1/2}\le C_{\rm pop}\overline{\mathcal J}(p),
 \quad \log\overline T\le18\nu^2,\quad
 \log\overline Q\le15+2\nu\le23\nu/4.
\]
The superscript \(1/2\) labels \(\theta\).
For every \(L\ge64\nu\ge256\),
\begin{equation}
 \log(L+4)\le L/32,\quad p+1\le129L/256,\quad
 2p+4\le65L/64.
 \label{eq:linear-elementary-bounds}
\end{equation}
Indeed \(\e^8>260\) by the exponential series through degree four,
and \(L/32-\log(L+4)\) increases on this interval.
Since \(\log24+(\nu-1)\log3\le\nu^2\),
the local logarithmic cost satisfies
\begin{align}
 \Lp(8C_{\mathrm o}Q_p)
 &\le37\nu^2+\log(p+2)+(p+1)\log\overline Q
                         +\log((2p+4)!)\nonumber\\
 &\le\frac{592+8+2967+2080}{65536}L^2
   <\frac3{32}L^2.
 \label{eq:linear-local-cost}
\end{align}
Here \(m!\le m^m\) pays the factorial, and the four rational
summands bound the four terms in the first line, respectively.
Also \(8C_Q\ge392C_{\rm pop}\e>392C_{\rm pop}\), so the same
bound applies to \(\Lp(392C_{\mathrm o}A_{\rm art,p}^{1/2})\).
The unchanged floor lower bound gives
\[
 (p+3/4)\log X-3L^2/8\ge7L^2/64>3L^2/32,
\]
paying the local half-target gate with strict slack.

For the true tails, the preceding moment estimates give
\[
 \Lp(392C_{\mathrm o}A_{\rm mod}^{1/2}),\quad
 \Lp(392C_{\mathrm o}A_{\rm far})
 \le18+2\nu+18\nu^2\le20\nu^2<L^2/32.
\]
Since \(\log X<L+3\), the moderate height power costs at most
\(L^2/4\), and the artificial one at most \(L^2/16\).
Their combined right sides are below \(9L^2/32\) and \(5L^2/32\),
respectively, both smaller than \(61L^2/8\).
Finally \(X>24L\), \(X>64\pi\), whence
\[
 \pi X/8<X^2/512,\qquad
 X^2/512>9L^2/8>(3/8+1/32)L^2.
\]
This pays the remote half-target gate. The factors \(8,392\),
and hence the half-target reservation for evaluation, are unchanged.

For completeness every non-rate gate holds on this smaller domain.
Use \(W\le\nu-1<L/64\), \(V\le12L\) and
\(X>12\e^L\ge2L^3\).
These give \(X>X_{\rm cl}=64\nu\), \(X>4(W+1)\),
\(V<X/4\), and \(2(V+2)<X/4\).
The strong ellipse and cutoff-span estimates already proved for
\(L\ge16\) apply. Also
\[
 K\le L\le(X-V)/(4\pi),\quad
 DV/X<24/L^2<1/2,\quad 2AV^2/X<128/L\le1/2,
\]
and \(V>\max(2R,W+1,64)\).
Since \(X>30\) and \(\pi<16/5\), the stronger slope gate holds.
The separate damping gate follows from
\[
 2\{(3/X)(\pi/4+L/2)+9/X^2\}
 <3(L+2)/X+18/X^2<1,
\]
using \(X>6(L+2)\) and \(X>6\).
Finally \(C_{\rm cal}=\nu+31<L\), so
\(C_{\rm cal}(1+L)<2L^2<X\).
Strong geometry supplies the ellipse logarithm bound as before.
Provider tolerances and strict descriptor witnesses retain their
separate finite contracts.

All estimates hold continuously for every \(L\ge64\nu\), using
\(L/2-1\le p\le L/2\), and the full-population bounds and exact cutoff
transfer cover every later \(N\) jump. No monotonicity of actual
factorial costs, or differentiation across chart changes, is asserted.
Lastly \(\log2>1/2\), \(\log(4\pi)<3\) give
\(\log(X_{\rm new}/(4\pi))>64\nu+1\).
\end{proof}

The old and new sufficient log-heights have exact ratio
\(L_*/L_{\rm new}=2^{16}\nu^3\ge2^{22}\).
For example on the dyadic chart family \(r=b=(\nu-3)/2\),
\(\nu=4,8,16\) gives new bounds \(256,512,1024\), against
\(2^{30},2^{34},2^{38}\).
These are symbolic comparisons of sufficient bounds, not field
measurements. The proof retains every arithmetic mode and completion
coefficient, both arms, the bounded integral, all buffers and channels,
and the fixed time interval. \Cref{sec:moving-endpoint} separately proves
\([1/L,2]\) at sufficient \(1024\nu^2\), with its own kernel and delivery
payments. Optimal onset, useful runtime, strong precision through zero
and complete evaluator certification remain separate questions.
